\chapter{Discussion and Conclusions}\label{ch:conclusions}

\section{Summary of findings, by research question}

\paragraph{RQ1 --- the exact joint score and its lifecycle.} For the
stationary AR(1) chain under coordinatewise OU corruption, the joint score
is the linear field $\score(x,t) = -\Qt x$ with
$\Qt = (e^{-2t}\Sigz + \Deltat I)^{-1}$, equal --- by Tweedie's identity,
for any prior --- to the amplified posterior-mean form
$(e^{-t}\E[a|x] - x)/\Deltat$. Its coupling structure is a complete,
closed-form lifecycle: exactly tridiagonal at $t = 0$ (the
Hammersley--Clifford fingerprint of Markovianity), band-filling at rate
$(\Qt)_{i,i+d} = (-1)^{d-1}(2t)^{d-1}(\Qz^d)_{i,i+d} + O(t^d)$ --- with a
transcription error of an earlier note corrected and the correction
audited --- and returning to the identity at the universal rate $e^{-2t}$.
In the bulk, one dimensionless dial $q(\alpha, t) \in (0, |\alpha|]$
governs simultaneously the posterior correlation length and the geometric
decay of evidence influence, with the exactly right limits $q \to 0$
($t \to 0$) and $q \to |\alpha|$ ($t \to \infty$).

\paragraph{RQ2 --- local computation of a global object.} The posterior of
any Markov chain observed through per-frame noise lives on a tree-shaped
factor graph, so sum--product message passing computes every marginal ---
hence every score component --- exactly, in $O(K)$. On the Gaussian chain
the messages remain Gaussian by \emph{algebraic closure} (two lemmas:
products and affine marginalisation), not by any central-limit argument;
the resulting algorithm is, update for update, the Kalman filter plus the
Rauch--Tung--Striebel smoother, and it reproduces the matrix score to
$10^{-14}$. The joint diffusion score of a Markov sequence \emph{is} a
smoothing problem --- a bridge that imports seventy years of estimation
theory into generative modelling.

\paragraph{RQ3 --- the price of the Gaussian closure.} The answer splits
into three cleanly separated statements. On the Gaussian chain, the
\emph{score} is closure-independent: mean field, BP, and AMP all solve the
same linear system for the posterior mean, so all return the exact score.
The \emph{variance} closure is where AMP pays: the per-node (TAP) closure
differs from the exact cavity by a single factor of two in the
discriminant, which produces an exact breakdown time $t_c(\alpha)$ and a
critical coupling $\alpha_c = \sqrt{2} - 1$ --- a sharp, closed-form
verdict on the ``central limit theorem on two data points''. Beyond
Gaussianity, on Laplace and other non-Gaussian chains measured against
validated grid quadrature, the Gaussian message closure costs tens of
percent of relative score error at small diffusion time (median $0.39$ at
$t = 0.02$; worst for bimodal innovations), decaying below $10^{-3}$ once
the diffusion noise Gaussianises the posteriors --- with the error
amplified at small $t$ by the exact factor $e^{-t}/\Deltat \sim 1/2t$, and
mitigated, counterintuitively, by \emph{stronger} chain correlation.

\paragraph{RQ4 --- implications for learned scores.} Locality has a price
and the price has a formula. The influence of evidence at distance $d$
decays as $q(\alpha,t)^d$; the error of a radius-$r$ local estimator
decays as $q^r$; and the penalty peaks in the mid-diffusion window, where
the band of $\Qt$ is fullest --- precisely where generative sampling
spends its critical steps. At equal locality budget, truncating the
\emph{inference} (windowed message passing) beats truncating the
\emph{matrix} (banding the score): $12.3\%$ against $21.0\%$ relative RMS
at $t = 0.05$. For architecture design the readings are direct: local,
convolution-like score heads are near-optimal at small and large $t$; the
receptive field should be sized as $r \gtrsim \xi \log(1/\varepsilon)$
with $\xi = 1/\log(1/q)$, or grown with diffusion time; and if a local
computation must be used, it should be a local \emph{algorithm}, not a
banded linear readout.

\section{Contributions to the literature}

To the exactly solvable programme in the statistical physics of diffusion
\citep{biroli2023generative, biroli2024dynamical, ghio2024sampling}, this
thesis adds the missing axis of \emph{internal structure}: a solvable
model in which a single data point has a dynamics of its own, and in which
the interplay of internal correlation ($\alpha$) and diffusion time ($t$)
is charted exactly. To the message-passing literature
\citep{mezard2009information, zdeborova2016statistical}, it contributes a
minimal, fully worked case where BP and AMP are \emph{both} solvable and
\emph{provably different} --- the factor-two discriminant, the phase
boundary $t_c(\alpha)$, $\alpha_c = \sqrt{2}-1$ --- together with a
quantitative audit of the Gaussian closure on chains, the regime where its
classical CLT justification fails by construction. To the practice of
diffusion modelling, it offers exact baselines: any learned score on
Markov-like sequence data can now be compared, regime by regime, against
a closed-form ground truth with known locality laws.

\section{Limitations}

Honesty about scope is part of the method. (i) The exact solutions concern
\emph{one-dimensional} state per frame and \emph{linear} (AR(1)) dynamics;
richer per-frame state keeps the tree structure but complicates the
algebra. (ii) The non-Gaussian results are numerical, not closed-form:
grid BP is an exact reference only within its validated resolution regime
($\sqrt{2t} \gtrsim 3\,\dd x$), and the Laplace closure error is a
measured curve, not a formula. (iii) The locality laws are proved for the
Gaussian bulk; their non-Gaussian counterparts are only structurally
delimited (Chapter~\ref{ch:laplace}). (iv) No learning is performed: the
architecture implications are readings of exact inference results, not
training experiments. (v) An earlier conjecture on low-rank structure of
the Laplace $K = 2$ Hessian proved regime-dependent and is deliberately
not asserted. (vi) Reverse-generation dynamics --- how score errors
integrate along the sampling trajectory --- was explored only
preliminarily and is not claimed here.

\section{Future work}

Four directions follow naturally, in increasing ambition.
\emph{Mixture messages.} The single-Gaussian closure fails worst on
bimodal beliefs; a two-component mixture closure on the same sweeps would
separate the message-representation error from the model-mismatch error,
and is the clearest next algorithmic step.
\emph{Discrete chains.} Replacing the continuous state with a finite
alphabet (transition-matrix dynamics, OU noise only in the channel) makes
BP messages finite vectors --- exact inference with no closure at all ---
and connects directly to existing tree-reconstruction theory; this was the
alternative route sketched in supervision.
\emph{Approximate Markovianity.} Real sequences are only approximately
Markov; preliminary experiments with a chain-plus-global-latent model
(rank-one residual) and with hybrid BP-plus-learned-residual scores
suggest the BP baseline absorbs most of the structure at a fraction of the
parameters, and deserve a systematic treatment.
\emph{Reverse dynamics.} Integrate the reverse SDE under exact,
Gaussian-closed, and truncated scores, and measure where along the
trajectory the closure and locality errors actually matter for sample
quality --- the pointwise-versus-dynamical error distinction.

\section{Closing remark}

The thesis began with a black box --- a neural network asked to learn
$\nabla \log p_t$ --- and ended with a dial, $q(\alpha, t)$, and a phase
boundary, $\alpha_c = \sqrt{2} - 1$. The distance between those two
endpoints is the argument for exactly solvable models: when the data have
structure, the score inherits it in a computable way, and both the
algorithms that exploit it (message passing) and the shortcuts that betray
it (closures, truncations) can be judged against ground truth rather than
against each other. For sequential data with Markov structure, this
programme is now complete at the Gaussian core and quantitatively mapped
at its non-Gaussian rim; the frontier --- mixtures, discreteness,
approximate Markovianity, dynamics --- is stated, and it is sharp.
